\documentclass[onecolumn, draftclsnofoot,10pt, compsoc]{IEEEtran}
\usepackage{graphicx}
\usepackage{url}
\usepackage{setspace}
\usepackage{listings}

\usepackage{geometry}
\geometry{textheight=9.5in, textwidth=7in}

% 1. Fill in these details
\def \CapstoneTeamName{Manovacumeter Team}
\def \CapstoneTeamNumber{20}
\def \GroupMemberOne{Cade Raichart}
\def \GroupMemberTwo{Daniel Kato}
\def \GroupMemberThree{Nathan Shepherd}
\def \CapstoneProjectName{ISS Barometer App }
\def \CapstoneSponsorCompany{NASA}
\def \CapstoneSponsorPerson{Don Pettit}

% 2. Uncomment the appropriate line below so that the document type works
\def \DocType{Progress Report}

\newcommand{\NameSigPair}[1]{\par
\makebox[2.75in][r]{#1} \hfil 	\makebox[3.25in]{\makebox[2.25in]{\hrulefill} \hfill		\makebox[.75in]{\hrulefill}}
\par\vspace{-12pt} \textit{\tiny\noindent
\makebox[2.75in]{} \hfil		\makebox[3.25in]{\makebox[2.25in][r]{Signature} \hfill	\makebox[.75in][r]{Date}}}}
% 3. If the document is not to be signed, uncomment the RENEWcommand below
%\renewcommand{\NameSigPair}[1]{#1}

%%%%%%%%%%%%%%%%%%%%%%%%%%%%%%%%%%%%%%%
\begin{document}
\begin{titlepage}
    \pagenumbering{gobble}
    \begin{singlespace}
    	\includegraphics[height=4cm]{coe_v_spot1}
        \hfill
        % 4. If you have a logo, use this includegraphics command to put it on the coversheet.
        \includegraphics[height=4cm]{NASA}
        \par\vspace{.2in}
        \centering
        \scshape{
            \huge CS Capstone \DocType \par
            {\large\today}\par
            \vspace{.5in}
            \textbf{\Huge\CapstoneProjectName}\par
            \vfill
            {\large Prepared for}\par
            \Huge \CapstoneSponsorCompany\par
            \vspace{5pt}
            {\Large\NameSigPair{\CapstoneSponsorPerson}\par}
            {\large Prepared by }\par
            Group\CapstoneTeamNumber\par
            % 5. comment out the line below this one if you do not wish to name your team
            %\CapstoneTeamName\par
            \vspace{5pt}
            {\Large
                \NameSigPair{\GroupMemberOne}\par
                \NameSigPair{\GroupMemberTwo}\par
                \NameSigPair{\GroupMemberThree}\par
            }
            \vspace{20pt}
        }
        \begin{abstract}
        % 6. Fill in your abstract
        The International Space Station, the conglomerate work of multiple national space associations, including NASA, requires precision devices to monitor the many aspects of space life.
        One very important aspect for the success of life in near vacuum is sustained air pressure inside the modules.
        This air pressure is monitored by mechanical barometers, called Manovacometers, which must be visually monitored, and are no longer available.
        The \CapstoneProjectName seeks to provide a suitable replacement for the current barometer device aboard the International Space Station.
        This document will serve to introduce the \CapstoneProjectName, and will detail the problem, the solution, and the success metrics of the project.
        The sections include the Definition of the Problem, Definition of Project, and the Metrics of Success.
        \end{abstract}
    \end{singlespace}
\end{titlepage}
\newpage
\pagenumbering{arabic}
\tableofcontents
% 7. uncomment this (if applicable). Consider adding a page break.
%\listoffigures
%\listoftables
\clearpage

% 8. now you write!
\section{Project Purpose}

\section{Project Goals}

\section{Current State of the Project}
% If anyone can think of a better way to state this, feel free to change the heading

\section{Project Problems}

\section{Interesting Code}
\begin{lstlisting}
  func startDisplayingPressureData() {
      self.altimeter.startRelativeAltitudeUpdates(to: OperationQueue.main, withHandler: {
        data, error in
          self.handlePressureReading(data: data!)
      })
  }

\end{lstlisting}

One of the main features of our application is its ability to read data from the iPad's built in barometer and display the data collected in different ways.
The way that data is collected from the barometer is quite interesting.
The process of reading data from the barometer requires a bit of processing power and time which can block other time sensitive tasks form being performed.
As a result, the measurement is done asynchronously through an object called the \texttt{OperationQueue}.
A method provided by the CoreMotion library called \texttt{startRelativeUpdates()} takes the \texttt{OperationQueue}'s \texttt{main} queue, and a callback function as arguement.
The function then appends a task to the queue to perform in the background, and when the task is complete, the result is passed into the callback function for further processing.\\

\section{Document Summaries}

\subsection{Problem Statement}

\subsection{Requirements Document}

\subsection{Technology Review}

\subsection{Design Document}

\section{Retrospective}
\begin{center}
	\begin{tabular}{| l | p{0.3\linewidth} | p{0.3\linewidth} | p{0.3\linewidth} |}
		\hline
		Week 1 & Positives & Deltas & Actions \\ \hline
		Week 2 & & & \\ \hline				% Put content between '&'s
		Week 3 & Finished problem statements and turned into both our client and teacher
 & Figure out how to share app with client, whom does not know the github system
 & step by step instructions on how to download the app \\ \hline
		Week 4 & & & \\ \hline
		Week 5 & Started Requirements Document rough draft & Need to start Requirements Document & Requirements Document will be started \\ \hline
		Week 6 & First prototype with two basic display of barometer numbers
  finished the Requirements Document & & \\ \hline
		Week 7 & Started Technology Review & Need to Start Technology Review & Technology Review will be started \\ \hline
		Week 8 & Completed Technology Review document rough draft & Need to complete Technology Review rough draft & A rough draft of the Technology Review will be created \\ \hline
		Week 9 &Finished Technology document
Started working on design document
 & Could not find IEEE paper & Searched more and asked around\\ \hline
		Week 10 & Completed Design Document and Progress Report & Need to complete Design Document and Progress Report & Design Document and Progress Report will be completed \\ \hline

	\end{tabular}
\end{center}




\end{document}
